% =====================================================================
% Appendix E --- Matrix variants and gamma/beta representation robustness
% ---------------------------------------------------------------------
% Job: where the presentation-decision variants live; all tables cite run
%      ids. Tally: fifteen matrix cells + twelve ladder cells + six
%      two-group cells = 33.
% Mindmap: N11, N12.
% Status: STUB --- narrative outline v3 (2026-09-20).
% =====================================================================
\section{Matrix variants and gamma/beta representation robustness}\label{app:variants}

% TABLE (linked, revision/tables/matrix_5x3_v10_flows.md): the full 15-cell
% matrix + the 18 sectoral cells, all run ids.
% ADDED (2026-10-18): the pricing-door profiles (tableE_pricing-profiles.tex),
% moved here from Section 6.4 when the main text switched to the goods-only
% profile.
% TODO: the full 15-cell matrix (citing matrix_5x3_v10 run ids); the
% gamma/beta variant tables; the ADR-0021 wage-structure axis as a
% representation-robustness row; the 33-cell tally
% (15 matrix + 12 ladder + 6 two-group).

\subsection{The pricing-door profiles and the mobile robustness}\label{app:profiles}

The main text carries the pricing door in the profile in which the utilisation
elasticity is confined to the goods sectors, because the entire direct incidence
of the programme falls there and because \textcite{Kaldor1966} reads Verdoorn's
law as a law of manufacturing. Table~\ref{tab:pricing-profiles} prints the rows
this leaves out: the counterfactual in which $\delta$ extends to every sector,
and both doors under mobile labour with the CPI numeraire.

The all-sector profile is worth stating explicitly, because it is what makes the
sign of the price response look like a property of the financing closure. With
$\delta$ on 36 service sectors, which receive none of the programme directly,
Kaldor--Verdoorn lowers the GDP deflator only under external finance
(\texttt{pricing\_capacity-v2-GAMMA-\{F1,F2,F3\}-dm05}: $+0.14$, $+0.03$,
$-0.82$ percent), against $-0.19$ to $-0.35$ percent in all three columns when
the elasticity is confined to goods. The reversal is produced by the service
sectors that lose demand, and with it productivity, as the programme reallocates
activity towards construction and its suppliers --- not by the sectors the
programme funds. It is a statement about where returns are assumed to live.

The mobile-labour rows repeat both doors under the CPI numeraire: employment is
fixed there and prices move relative to the CPI, and F2 and F3 remain
indistinguishable in every row, so the financing neutrality of
Section~\ref{sec:gauge} survives the pricing door.

% =====================================================================
% Table E --- pricing-door profiles not carried in the main text (GENERATED)
% ---------------------------------------------------------------------
% Do NOT edit by hand. Regenerate with:  make_pricing_door.py
% OpenCGE tag v0.5.0 (ADR-0044); source run ids are listed below.
% Sources: pricing_capacity_rows-v2-<row>-<fin>-<dp05|dm05>,
% pricing_capacity-v3-GAMMA-<fin>-<dp05|dm05>,
% pricing_mobile-v3-ALPHA-<fin>-<dp05|dm05>.
% One panel: the all-sector capacity counterfactual; the kappa = 0.25
% distributed-profit rows were dropped from the paper on 2026-10-24
% (they repeated the invariance Section 6.2 already states).
% =====================================================================
\begin{table}[htbp]
\centering
\scriptsize
\setlength{\tabcolsep}{1.6pt}
\renewcommand{\arraystretch}{1.15}
\resizebox{\textwidth}{!}{%
\begin{tabular}{@{}l@{\hspace{6pt}}rrrrr@{\hspace{6pt}}rrrrr@{\hspace{6pt}}rrrrr@{}}
\toprule
 & \multicolumn{5}{c}{\textbf{F1}} & \multicolumn{5}{c}{\textbf{F2}} & \multicolumn{5}{c}{\textbf{F3}} \\
\cmidrule(lr){2-6}\cmidrule(lr){7-11}\cmidrule(lr){12-16}
\multicolumn{1}{c}{Labor closure and door} & \multicolumn{1}{c}{$\Delta P$} & \multicolumn{1}{c}{\shortstack{$\max$\\$|\Delta p|$}} & \multicolumn{1}{c}{$\Delta L$} & \multicolumn{1}{c}{$\Delta W_C$} & \multicolumn{1}{c}{$\Delta\Pi/Y$} & \multicolumn{1}{c}{$\Delta P$} & \multicolumn{1}{c}{\shortstack{$\max$\\$|\Delta p|$}} & \multicolumn{1}{c}{$\Delta L$} & \multicolumn{1}{c}{$\Delta W_C$} & \multicolumn{1}{c}{$\Delta\Pi/Y$} & \multicolumn{1}{c}{$\Delta P$} & \multicolumn{1}{c}{\shortstack{$\max$\\$|\Delta p|$}} & \multicolumn{1}{c}{$\Delta L$} & \multicolumn{1}{c}{$\Delta W_C$} & \multicolumn{1}{c}{$\Delta\Pi/Y$} \\
\midrule
\multicolumn{16}{@{}l}{\textit{Capacity effects on all sectors: No-reallocation}} \\
\quad $\delta = +0.5$ & $+0.74$ & $26.82$ & $0.00$ & $+0.10$ & $0.00$ & $+0.86$ & $35.47$ & $0.00$ & $-1.99$ & $0.00$ & $+0.86$ & $35.47$ & $0.00$ & $-1.99$ & $0.00$ \\
\quad $\delta = -0.5$ & $+0.48$ & $16.39$ & $0.00$ & $+0.06$ & $0.00$ & $+0.53$ & $19.52$ & $0.00$ & $-1.99$ & $0.00$ & $+0.53$ & $19.52$ & $0.00$ & $-1.99$ & $0.00$ \\
\midrule
\multicolumn{16}{@{}l}{\textit{Capacity effects on all sectors: Standard CGE}} \\
\quad $\delta = +0.5$ & $+0.12$ & $4.06$ & $0.00$ & $+0.15$ & $0.00$ & $+0.13$ & $4.27$ & $0.00$ & $-1.83$ & $0.00$ & $+0.13$ & $4.27$ & $0.00$ & $-1.83$ & $0.00$ \\
\quad $\delta = -0.5$ & $-0.16$ & $4.56$ & $0.00$ & $+0.07$ & $0.00$ & $-0.16$ & $4.47$ & $0.00$ & $-1.87$ & $0.00$ & $-0.16$ & $4.47$ & $0.00$ & $-1.87$ & $0.00$ \\
\midrule
\multicolumn{16}{@{}l}{\textit{Capacity effects on all sectors: Elastic labor supply ($\eta_s = 0.5$)}} \\
\quad $\delta = +0.5$ & $+0.38$ & $14.11$ & $+0.15$ & $+0.38$ & $0.00$ & $+0.42$ & $16.21$ & $+0.16$ & $-1.64$ & $0.00$ & $+0.42$ & $16.21$ & $+0.16$ & $-1.64$ & $0.00$ \\
\quad $\delta = -0.5$ & $+0.13$ & $4.31$ & $+0.07$ & $+0.40$ & $0.00$ & $+0.13$ & $4.54$ & $+0.08$ & $-1.57$ & $0.00$ & $+0.13$ & $4.54$ & $+0.08$ & $-1.57$ & $0.00$ \\
\midrule
\multicolumn{16}{@{}l}{\textit{Capacity effects on all sectors: Employment adjustment}} \\
\quad $\delta = +0.5$ & $+0.12$ & $4.06$ & $+0.04$ & $+0.21$ & $0.00$ & $+0.10$ & $4.22$ & $+0.35$ & $-1.30$ & $0.00$ & $-0.10$ & $3.90$ & $+2.61$ & $+2.14$ & $0.00$ \\
\quad $\delta = -0.5$ & $-0.20$ & $4.60$ & $-0.15$ & $-0.53$ & $0.00$ & $-0.16$ & $4.47$ & $-0.01$ & $-1.89$ & $0.00$ & $+0.12$ & $4.03$ & $+1.05$ & $+2.43$ & $0.00$ \\
\midrule
\multicolumn{16}{@{}l}{\textit{Capacity effects on all sectors: Leontief closure}} \\
\quad $\delta = +0.5$ & $+0.14$ & $4.52$ & $+0.02$ & $+0.21$ & $0.00$ & $+0.10$ & $4.54$ & $+0.31$ & $-1.34$ & $0.00$ & $-0.15$ & $4.21$ & $+2.45$ & $+1.89$ & $0.00$ \\
\quad $\delta = -0.5$ & $-0.16$ & $4.12$ & $-0.12$ & $-0.44$ & $0.00$ & $-0.13$ & $4.15$ & $+0.03$ & $-1.79$ & $0.00$ & $+0.14$ & $3.76$ & $+1.22$ & $+2.61$ & $0.00$ \\
\bottomrule
\end{tabular}}
\caption{The capacity door with the utilization elasticity extended to every sector ---
the all-sector counterfactual of Table~\ref{tab:pricing-kaldor} --- on all five
labor closures. Conventions as in Table~\ref{tab:pricing-markup}, except that the
welfare column is household consumption ($\Delta W_C$ of Table~\ref{tab:results-matrix}).}
\label{tab:pricing-profiles}
\end{table}
